\clearpage\fancyhead[L]{\small COMET.}\hypertarget{native-347-lower}{}\label{native-347-lower}
\begin{center}\textbf{COMET.}\par\vspace{4pt}\hrule\vspace{4pt}\fontsize{10}{12}\selectfont Göttingische gelehrte Anzeigen. No. 201, pp. 2001\,..\,2004, 19 December 1811.\par\vspace{4pt}\hrule\vspace{4pt}\end{center}

\par\indent In issue 151 of this journal [p. 337 of this volume] we reported the first observations made here of this year’s great comet, together with the first revised parabolic elements. Those elements were presented only as provisional. They soon began to deviate from the observations by a few minutes, and at the beginning of October the discrepancy had already grown to a quarter of a degree. A slight revision of the parabolic elements was enough to remove this error, however. This second set of revised elements was published in the October issue of the Monatl. Correspondenz [p. 338 of this volume]. By the beginning of November the error in these elements, which had been calculated from only a few individual observations, had again grown to a few minutes. Professor \textsc{Gauss} therefore thought it worth investigating to what extent this could already be regarded as proof that the orbit was elliptical. He had Mr. \textsc{Nicolai} carry out the necessary calculations under his supervision. We have already given examples in this journal of Mr. Nicolai’s outstanding skill and care in astronomical computation. The calculations were based on most of the observations Mr. von \textsc{Zach} made during the first period in which the comet was visible, and on a large number of more recent observations extending to 6 November. Four normal places were derived by comparing those observations with Professor \textsc{Gauss}’s latest elements. The result was that a very small correction to these latest elements sufficed to represent the more recent observations accurately, while only small differences remained for the earlier observations, namely $16''$ in longitude and $28''$ in latitude for the first normal place on 16 April, and $51''$ in longitude and $120''$ in latitude for the second on 18 May. There is no reason to suppose that the latter normal place, the mean of a large number of observations that were admittedly of limited precision, is actually wrong by the full amount of so large a discrepancy. Yet the discrepancy is still too small, given the uncertainty about how much of it is attributable to the normal place itself, to serve as a basis for a reasonably reliable determination of the ellipse. Professor Gauss therefore thought it better to postpone the work until later observations made it possible to decide the matter. In any case the orbital period is well over 1000 years. The revised parabolic elements according to Mr. \textsc{Nicolai}’s calculation are as follows.

\Needspace{18\baselineskip}
\begin{center}\fontsize{10}{12}\selectfont\setlength{\tabcolsep}{4pt}\renewcommand{\arraystretch}{1.15}\begin{tabular}{l}
Passage through perihelion, 12 September 1811\\
\quad $6^{\mathrm u}\;30^{\mathrm m}\;55^{\mathrm s}$ mean time in Göttingen\\
\end{tabular}\end{center}
\begin{center}\fontsize{10}{12}\selectfont\setlength{\tabcolsep}{3pt}\renewcommand{\arraystretch}{1.15}\begin{tabular}{llr}
Logarithm of the minimum distance &  & $0.0151048$\\
Longitude of perihelion & .\;.\;.\;.\;. & $75^\circ\;1'\;44''3$\\
Longitude of the ascending node & .\;.\;. & $140\;21\;57.5$\\
\end{tabular}\end{center}
\noindent Both are fixed relative to the stars and reckoned from the equinox of 12 September.
\begin{center}\fontsize{10}{12}\selectfont\setlength{\tabcolsep}{3pt}\renewcommand{\arraystretch}{1.15}\begin{tabular}{llr}
Inclination of the orbit & .\;.\;.\;.\;.\;.\;. & $73^\circ\;4'\;30''9$\\
\end{tabular}\end{center}
\noindent Motion retrograde.

\SourceLeaf{348}
\par\indent To facilitate observations in January of next year, Mr. \textsc{Nicolai} has also calculated the following ephemeris, which we give here in full, although the comet will cease to be visible sooner. \mbox{--- --- ---}


\par\indent The foregoing had already been written for publication when Professor \textsc{Gauss} received word in a letter from Mr. von \textsc{Lindenau} that Mr. \textsc{Pons} had discovered another new comet in Eridanus at Marseille on 16 November. The letter included the following observations made by Mr. von \textsc{Zach} at Marseille.

\begin{center}
\fontsize{10}{12}\selectfont\setlength{\tabcolsep}{3pt}\renewcommand{\arraystretch}{1.15}\begin{tabular}{l|c|c|c}
1811 & \shortstack{Mean time\\in Marseille} & \shortstack{Right\\ascension} & \shortstack{Southern\\declination} \\ \hline
Nov. 17 & $10^{\mathrm u}\;23^{\mathrm m}$ & $67^\circ.25'$ & $25^\circ\,52'$ \\
18 & $11\;\;11\;\;17^{\mathrm s}3$ & $67\;\;14\;\;39''8$ & $25\;\;24\;\;\;8''6$ \\
19 & $\;9\;\;59\;\;32.8$ & $67\;\;\;4\;\;59.6$ & $24\;\;54\;\;\;8.5$ \\
20 & $10\;\;\;8\;\;37.7$ & $.66\;\;56\;\;\;8.2$ & $24\;\;18\;\;\;9.2$ \\
21 & $10\;\;14\;\;45.5$ & $66\;\;46\;\;53.0$ & $23\;\;41\;\;47.8$ \\
\end{tabular}
\end{center}


\Needspace{22\baselineskip}\par\indent On the very evening of 9 December, when this news arrived, a very clear sky favored the search for the comet at our observatory. Professor \textsc{Harding} immediately sighted it with a comet finder near a star of the seventh magnitude. Professor \textsc{Gauss} compared the comet with that star several times using the ring micrometer. The star’s apparent position was determined from the \emph{Histoire céleste} to be $64^\circ\,11'\,23''1$ in right ascension and $10^\circ\,22'\,2''5$ in southern declination. From this the following position of the comet was obtained.

\Needspace{9\baselineskip}
\begin{quote}\raggedright 1811 Dec. 9. $10^{\mathrm u}\;6^{\mathrm m}\;52^{\mathrm s}$ R.A. $63^\circ\;49'\;41''4$, southern declination $10^\circ\;21'\;55''5$\par\medskip On 11 Dec., at $10^{\mathrm u}\;34^{\mathrm m}\;1^{\mathrm s}$ mean time, the position found was R.A. $63^\circ\;33'\;20''5$, S. decl. $8^\circ\;39'\;54''7$\end{quote}


\par\indent So far this new comet is very small and faint, about as bright as a star of the eighth magnitude. It seems, however, to be approaching the Earth, as its motion is accelerating. It can also be observed better with the ring micrometer than the great comet, since its center is more clearly defined.


\par\indent 12 December.


